\subsection{Explainable Predictors for Credit Risk Prediction}

    Incorporating explainability into credit risk prediction has been a growing area of study. These studies aim to achieve robust predictive performance while prioritising the explainability of predictions.\\

    A study by \cite{zhou2025enhancing} addresses credit risk prediction by comparing the predictive performance of XGBoost, RF, Least Squares SVM (LSSVM), and a CNN. After applying SMOTE to balance the dataset, XGBoost achieved the best predictive accuracy, reaching 92\%. To enhance predictor explainability, SHAP analysis was applied to quantify the individual contributions of each feature. From the predictors evaluated, XGBoost performed best in terms of predictive accuracy and also provided explanations through SHAP values.\\ 

    In \cite{baisholan2025fraudx}, the authors present an ensemble framework that combines RF and XGBoost using probability-weighted averaging. The framework was evaluated using eight predictors, including LR, DT, Gradient Boosting (GB), LightGBM (LGBM), MLP, Naive Bayes (NB), and XGBoost. Their proposed ensemble significantly outperformed all others, achieving an AUC of 97\%. SHAP was employed to explain the predictions, which allowed them to determine the contribution of each feature.\\

   In \cite{sun2025enhancing}, XGBoost is used as a predictor, combined with SHAP, LIME, and K-means to enhance explainability, achieving an AUC of 78\%. SHAP and LIME are used to generate feature-contribution vectors for each training instance, which are grouped using K-means. For explaining the prediction of a new instance, its feature-contribution vector is calculated and assigned to the nearest cluster, and the prediction explanation is provided as the cluster’s average vector.\\

    In \cite{nguyen2025comparative}, the authors present a comparative analysis of the performance of boosting algorithms, including LightGBM, XGBoost, CatBoost, and AdaBoost. Additionally, SMOTE-Tomek \cite{batista2003balancing} is employed to balance the dataset. The results indicate that LightGBM achieves the best performance with 95\%  accuracy and 98\% AUC, outperforming the other boosting methods. SHAP is applied to quantify the contribution of each feature to the prediction.\\

    In \cite{zandi2025attention}, a multilayer dynamic graph neural network (DYMGNN), proposed for credit risk prediction. The authors compare their approach against several baselines: LR, XGBoost, and DNN. Unlike many related studies, no resampling or class balancing methods are applied. Experimental results show that the DYMGNN achieves the best results with 81\% accuracy and 83\% F1-score. For explainability, SHAP is used to identify the features that contribute most to the predictions. \\

    In \cite{wang2025two}, the authors propose an explainable two-stage technique to explain the results of credit risk prediction. In the first stage, SHAP is applied to analyse the importance of features according to their positive or negative impact on the prediction. In the second stage, they use a genetic algorithm to generate counterfactual explanations by perturbing the features according to this categorisation, thereby identifying the minimum necessary changes to modify the prediction. This technique was evaluated using multiple predictors, including SVM, XGBoost, MLP, and LSTM, with an accuracy of 90\% for the employed predictors. The results demonstrate that the features modified most frequently closely align with the features identified as most relevant by SHAP.\\

    In \cite{aruleba2025enhanced}, the authors compare the performance of several Deep Learning (DL) architectures for credit risk prediction, including Convolutional Neural Networks (CNN), Long Short-Term Memory (LSTM) networks, Gated Recurrent Units (GRU), and Graph Neural Networks (GNN). To address class imbalance, they apply SMOTEENN. The experimental results show that GRU outperformed all other predictors, achieving an accuracy of 91\%. To explain the prediction, SHAP is used to quantify the importance of features and their contributions to the predictions.\\

    In \cite{hartomo2025novel}, the authors propose an approach for credit risk prediction on imbalanced datasets by integrating the TabTransformer model with a weighted loss function. This approach helps mitigate the bias toward the majority class, improving predictive performance on minority classes, achieving 91\% accuracy and an AUC of 93\%. SHAP is employed to evaluate feature importance and identify those features that contribute most to the predictions.\\
    
    In \cite{damanik2025advanced}, the authors propose a robust framework for credit risk prediction on imbalanced datasets by introducing the K-SMOTEENN method, which integrates oversampling and noise-filtering methods with a stacking ensemble. The ensemble combines Random Forest, XGBoost, and LightGBM, resulting in a predictive performance of 92\% F1-score, and 96\% AUC. SHAP is used to analyse and explain the contributions of individual features to the prediction.\\
    
    %In \cite{han2024nate}, the authors propose NATE, a non-parametric approach to credit risk prediction in imbalanced datasets that integrates GB as a predictor. The study compares GB with other classifiers, including LR, SVM, XGBoost, and RF, demonstrating that GB achieves the highest predictive performance with an AUC of 96\% and a F1-score of 90\%. NATE applies SMOTE to balance the dataset and employs SHAP to explain the prediction, highlighting the contribution of each feature to the prediction.\\ 

    %In \cite{talaat2024toward}, the authors propose CreditNetXAI, which combines deep neural networks with XAI techniques for credit risk prediction, using feature selection methods to identify the most relevant features for prediction. CreditNetXAI obtains an accuracy of 83\%. To explain predictions, they use SHAP to assign an importance to each feature and indicate its contribution to the prediction.\\ 

    %In \cite{zheng2024micro}, the authors employ Decision Tree-based LightGBM with an optimized gradient approach to credit risk prediction, demonstrating an accuracy of 93\% and an AUC of 80\%.  To explain the prediction, they use SHAP to show the importance of features in the prediction. Similarly in \cite{de2022explainable} used a LightGBM for credit risk prediction, achieving an accuracy of 82\% and SHAP to explain the prediction. In \cite{nallakaruppan2024explainable}, the authors propose using RF with explainable AI techniques for credit risk prediction, achieving an accuracy of 99\%. To generate explanations for predictions, LIME, SHAP, and Partial Dependence Plots (PDPs) were used, which help identify how feature values contribute to the prediction. It is important to note that these results were obtained on different datasets, making direct comparison between reported accuracies and AUCs not the same. \\

    As discussed in the above works, credit risk prediction has been addressed through a wide range of predictors, balancing methods, feature selection methods, and evaluated through multiple performance metrics. Nevertheless, accurately detecting defaults remains a significant challenge, as no single method has consistently outperformed others across different scenarios. Moreover, the most commonly used explainability approaches, such as SHAP and LIME, while helpful in assessing the individual contribution of features, fail to capture feature value combinations that characterise different types of applicants. This limitation is particularly relevant, since financial institutions increasingly demand not only high predictive accuracy but also clear and case-specific justifications, both to meet regulatory requirements and to improve communication with applicants. 